\clearpage
\gdef\otChapNum{0}\pagestyle{otFront}
\phantomsection\label{defis-debut}\pdfbookmark[0]{Défis d'approfondissement}{bm-defis}
\otPageTitre{Défis d'approfondissement}{Chercher plus loin avec les outils de Terminale}
\fontsize{11}{16}\selectfont
Les \otNbExosChapitres\ exercices des chapitres installent les connaissances et les méthodes.
Les \otNbDefis\ défis suivants, un par chapitre, portent le total à \textbf{\otNbExosTotal\ exercices corrigés}.
Ils prolongent chaque thème par une question qui demande de prendre du recul, de construire un modèle
ou de relier plusieurs notions.\par\vspace{4mm}
\textbf{Avancé} désigne une modélisation ou une synthèse soutenue.
\textbf{Très exigeant} ajoute une discussion, une démonstration ou plusieurs changements de point de vue.
Ces indications sont des repères pédagogiques, pas un classement officiel.\par\vspace{4mm}
Les défis sont facultatifs pour un premier passage. Ils restent dans le cadre du programme ; un outil qui
le dépasserait est fourni dans l'énoncé. Les durées proposées sont indicatives.\par\vspace{5mm}
\editionEtape{Chercher avant de lire la solution}
Faire un schéma, écrire les données et les lois utiles, tester un cas simple, puis noter ce qui manque
pour conclure. En cas de blocage, lire seulement le premier indice, puis reprendre la recherche.
Les indications sont séparées des énoncés et des corrigés.
\begin{center}\renewcommand{\arraystretch}{1.25}
\begin{tabular}{@{}lll@{}}\toprule
Défis & Thèmes & Préparation conseillée\\\midrule
D01 à D12 & Constitution et transformations de la matière & Chapitres 1 à 12\\
D13 à D17 & Mouvement et interactions & Chapitres 13 à 17\\
D18 à D20 & L'énergie : conversions et transferts & Chapitres 18 à 20\\
D21 à D26 & Ondes et signaux & Chapitres 21 à 26\\
D27 & Mesures et incertitudes & Chapitre 27\\\bottomrule
\end{tabular}\end{center}
\editionRetenir{Une solution incomplète mais justifiée est un progrès.
Après la correction, refaire la question clé sans regarder, puis expliquer pourquoi la méthode fonctionne.}

\editionPage{defi-ex-1}{D01 · Mélanger n'est pas moyenner}{Recherche et rédaction}{Très exigeant · 45 à 60 min · Chapitre 1\hfill\hyperref[defi-cor-1-1]{Corrigé p.~\pageref*{defi-cor-1-1}}}
\begin{blocEdition}{defi-ex-1}
Toutes les solutions considérées sont des solutions d'acide chlorhydrique : le chlorure d'hydrogène
y a totalement réagi avec l'eau. On note $h=\hho/\cst$. Sauf à la question 4, on néglige les ions
apportés par l'eau elle-même.
\label{figedition-defi-ex-1-1}\schema{defi-01-melange}
\begin{questions}
\item On mélange 100,0 mL d'une solution de pH 2,0 et 100,0 mL d'une solution de pH 4,0.
Calculer le pH du mélange. Est-il égal à la moyenne des deux pH ?
\item On mélange des volumes égaux de deux solutions de pH $p_1$ et $p_2$. Montrer que le pH $p$ du
mélange vérifie $p\leq\frac{p_1+p_2}{2}$, avec égalité si et seulement si $p_1=p_2$.
On pourra utiliser l'inégalité $\frac{a+b}{2}\geq\sqrt{ab}$, valable pour $a>0$ et $b>0$.
\item À 100,0 mL de la solution de pH 2,0, quel volume de la solution de pH 4,0 faut-il ajouter pour
obtenir un pH de 3,0 ? Comparer au volume d'eau pure qui produirait le même effet, et expliquer la différence.
\item On dilue une solution d'acide chlorhydrique jusqu'à $c=\qty{1,0e-8}{mol.L^{-1}}$.
Pourquoi la relation $\pH=-\log(c/\cst)$ conduit-elle à un résultat absurde ?
On admet qu'en solution aqueuse $\hho\times[\ce{HO-}]=K_\mathrm{e}\times{\cst}^2$, avec
$K_\mathrm{e}=\num{1.0e-14}$ à 25 °C, et que la solution est électriquement neutre :
$\hho=[\ce{Cl-}]+[\ce{HO-}]$. Établir l'équation $h^2-\frac{c}{\cst}\,h-K_\mathrm{e}=0$ et en déduire le pH.
\item Montrer que l'écart entre le pH exact et $-\log(c/\cst)$ est inférieur à 0,01 dès que
$c>\qty{6,5e-7}{mol.L^{-1}}$. Écrire une fonction Python \code{ph\_exact(c)} qui renvoie le pH exact.
\end{questions}
Les questions 1 à 3 n'utilisent que la définition du pH ; les questions 4 et 5 s'appuient sur les relations admises.
\end{blocEdition}

\editionPage{defi-ex-2}{D02 · La méthode des ajouts dosés}{Recherche et rédaction}{Très exigeant · 45 à 60 min · Chapitre 2\hfill\hyperref[defi-cor-2-1]{Corrigé p.~\pageref*{defi-cor-2-1}}}
\begin{blocEdition}{defi-ex-2}
Dans un échantillon complexe, comme un jus de fruit ou une eau de rivière, d'autres espèces peuvent modifier
l'absorbance de l'espèce dosée : une droite d'étalonnage préparée dans l'eau pure ne convient plus.
La méthode des ajouts dosés contourne cette difficulté.\par\smallskip
Dans cinq fioles jaugées de 50,0 mL, on introduit $V_0=\qty{10,0}{mL}$ de l'échantillon, de concentration inconnue
$c_0$ en espèce colorée X, puis un volume $V_\mathrm{s}$ d'une solution étalon de X de concentration
$c_\mathrm{s}=\qty{1,00e-4}{mol.L^{-1}}$ ; on complète avec de l'eau distillée. On mesure l'absorbance $A$
de chaque solution, dans la même cuve et à la même longueur d'onde.
\begin{center}\small\renewcommand{\arraystretch}{1.2}
\begin{tabular}{@{}l*{5}{c}@{}}\toprule
$V_\mathrm{s}$ (mL) & 0 & 2,0 & 4,0 & 6,0 & 8,0\\
$A$ & 0,209 & 0,311 & 0,417 & 0,520 & 0,623\\\bottomrule
\end{tabular}\end{center}
\begin{questions}
\item Pourquoi une droite d'étalonnage réalisée dans l'eau pure peut-elle conduire à un résultat faux ?
\item On admet que le coefficient $k$ de la loi $A=k\,c$ est le même dans les cinq fioles. Montrer que
$A=\dfrac{k}{V_\mathrm{T}}\left(c_0V_0+c_\mathrm{s}V_\mathrm{s}\right)$, avec $V_\mathrm{T}=\qty{50,0}{mL}$.
\item En déduire que $A$ est une fonction affine de $V_\mathrm{s}$. Exprimer l'abscisse du point d'intersection de
la droite avec l'axe des volumes en fonction de $c_0$, $V_0$ et $c_\mathrm{s}$.
\item Modéliser les mesures par une fonction affine, à la calculatrice ou en Python, puis déterminer $c_0$ et $k$.
\item La matrice ajoute en réalité une absorbance parasite constante $A_\mathrm{p}=0{,}020$, présente dans les cinq
fioles et non retirée. Quelle valeur de $c_0$ aurait-on dû trouver ? Quelle est l'erreur relative commise ?
Proposer un moyen de l'éviter.
\end{questions}
Les volumes de l'échantillon et de la fiole sont les mêmes pour les cinq mesures : seule la quantité ajoutée varie.
\end{blocEdition}

\editionPage{defi-ex-3}{D03 · Identifier un acide inconnu}{Recherche et rédaction}{Très exigeant · 50 à 70 min · Chapitres 1 à 3\hfill\hyperref[defi-cor-3-1]{Corrigé p.~\pageref*{defi-cor-3-1}}}
\begin{blocEdition}{defi-ex-3}
Un flacon contient un solide blanc, soluble dans l'eau, dont l'étiquette a disparu : on sait seulement qu'il s'agit
d'un acide carboxylique. On réalise quatre analyses.
\begin{questions}[label=\otLabel{\Alph*.},leftmargin=1.9em]
\item Titrage : 0,150 g de solide, dissous dans 50,0 mL d'eau, sont titrés par la soude à \qty{0,100}{mol.L^{-1}}.
La courbe pH-métrique ne présente qu'un seul saut de pH, pour $V_\mathrm{E}=\qty{25,4}{mL}$.
\item Analyse élémentaire, en pourcentages massiques : C : 40,7 \% ; H : 5,1 \% ; O : 54,2 \%.
\item Spectre infrarouge : bande très large entre 2\,500 et 3\,200 cm$^{-1}$ et bande forte vers 1\,700 cm$^{-1}$.
\item Température de fusion mesurée au banc Kofler : 184 °C.
\end{questions}
\begin{center}\small\renewcommand{\arraystretch}{1.02}
\begin{tabular}{@{}lll@{}}\toprule
Acide & Formule semi-développée & Température de fusion\\\midrule
éthanoïque & \ce{CH3-COOH} & 17 °C\\
benzoïque & \ce{C6H5-COOH} & 122 °C\\
méthylpropanedioïque & \ce{CH3-CH(COOH)2} & 135 °C\\
butanedioïque & \ce{HOOC-CH2-CH2-COOH} & 185 °C\\\bottomrule
\end{tabular}\end{center}
\begin{questions}
\item Que révèle le spectre infrarouge ?
\item Calculer la quantité d'ions hydroxyde versée à l'équivalence. Quelle serait la masse molaire de l'acide s'il ne
possédait qu'un groupe carboxyle ? Montrer que l'analyse élémentaire conduit alors à une formule \ce{C_xH_yO_z}
impossible : dans une telle molécule, le nombre d'atomes d'hydrogène est toujours pair.
\item L'acide possède $p$ groupes carboxyle. Exprimer sa masse molaire en fonction de $p$. On admet qu'une molécule
\ce{C_xH_yO_z} sans cycle compte $\frac{2x+2-y}{2}$ doubles liaisons ; en déduire que $p=2$ et donner la formule brute.
\item Identifier l'acide. Pourquoi l'analyse élémentaire seule ne permettait-elle pas de conclure ?
\item Pourquoi n'observe-t-on qu'un seul saut de pH, alors que la molécule possède deux groupes carboxyle ?
On admet que les deux acidités de la molécule sont voisines.
\end{questions}
\end{blocEdition}

\editionPage{defi-ex-4}{D04 · Ordre 1 ou ordre 2 ?}{Recherche et rédaction}{Avancé · 40 à 50 min · Chapitre 4\hfill\hyperref[defi-cor-4-1]{Corrigé p.~\pageref*{defi-cor-4-1}}}
\begin{blocEdition}{defi-ex-4}
Toutes les transformations ne suivent pas une loi de vitesse d'ordre 1. On étudie à 25 °C la saponification de
l'éthanoate d'éthyle par la soude :
\[\ce{CH3COOC2H5 + HO- -> CH3COO- + C2H5OH}.\]
Les deux réactifs sont introduits à la même concentration $c_0=\qty{1,00e-2}{mol.L^{-1}}$ : leurs concentrations
restent égales à chaque instant ; on les note $c$. Un suivi conductimétrique donne :
\begin{center}\small\renewcommand{\arraystretch}{1.2}
\begin{tabular}{@{}l*{7}{c}@{}}\toprule
$t$ (s) & 0 & 120 & 300 & 600 & 900 & 1\,200 & 1\,800\\
$c$ (mmol$\cdot$L$^{-1}$) & 10,00 & 8,83 & 7,52 & 6,02 & 5,03 & 4,31 & 3,36\\\bottomrule
\end{tabular}\end{center}
\begin{questions}
\item Estimer le temps de demi-réaction à l'aide du tableau. Si la loi de vitesse était d'ordre 1, quelle
concentration devrait-on mesurer à $t=\qty{1800}{s}$ ? Conclure.
\item Calculer $\ln(c/c^\circ)$, avec $c^\circ=\qty{1}{mol.L^{-1}}$, à chaque date, puis les coefficients directeurs
des segments qui joignent deux points successifs du graphique $\ln(c/c^\circ)=f(t)$. Pourquoi ce calcul
confirme-t-il la conclusion précédente ?
\item On admet que, pour cette réaction, la vitesse volumique de disparition vaut $v=k\,c^2$ (loi de vitesse
d'ordre 2) et qu'alors $\dfrac{1}{c}=\dfrac{1}{c_0}+k\,t$. Tester ce modèle, à la calculatrice ou en Python, et
déterminer la constante de vitesse $k$.
\item Exprimer le temps de demi-réaction en fonction de $k$ et de $c_0$, le calculer et le comparer au résultat de la
question 1. Que deviendrait-il pour $c_0=\qty{5,0e-3}{mol.L^{-1}}$ ?
\item Montrer que la durée nécessaire pour diviser à nouveau la concentration par deux, de $c_0/2$ à $c_0/4$, est
double de la précédente. À quelle date la concentration vaut-elle $c_0/4$ ? Comparer avec une loi d'ordre 1 de même
temps de demi-réaction.
\end{questions}
Les lois de vitesse d'ordre 2 ne sont pas au programme : la relation utile est fournie dans l'énoncé.
\end{blocEdition}

\editionPage{defi-ex-5}{D05 · Une horloge à iode}{Recherche et rédaction}{Très exigeant · 50 à 60 min · Chapitres 4 et 5\hfill\hyperref[defi-cor-5-1]{Corrigé p.~\pageref*{defi-cor-5-1}}}
\begin{blocEdition}{defi-ex-5}
Dans une « horloge à iode », un mélange incolore devient brusquement bleu nuit au bout d'une durée $t_\mathrm{c}$
très reproductible. Il contient de l'eau oxygénée, des ions iodure, un excès d'acide, de l'empois d'amidon et un peu
d'ions thiosulfate \ce{S2O3^2-}. Deux réactions se produisent :
\[\text{(1), lente : }\ce{H2O2 + 2 I- + 2 H3O+ -> I2 + 4 H2O}\qquad\text{(2), très rapide : }\ce{I2 + 2 S2O3^2- -> 2 I- + S4O6^2-}\]
Tant qu'il reste des ions thiosulfate, le diiode formé par (1) est aussitôt consommé par (2) ; la couleur bleue du
diiode en présence d'amidon apparaît dès que les ions thiosulfate sont épuisés. Dans tous les mélanges,
$[\ce{S2O3^2-}]_0=\qty{5,0e-4}{mol.L^{-1}}$ et les ions oxonium sont en large excès.
\begin{center}\small\renewcommand{\arraystretch}{1.2}
\begin{tabular}{@{}l*{5}{c}@{}}\toprule
Expérience & 1 & 2 & 3 & 4 & 5\\\midrule
$[\ce{H2O2}]_0$ (mol$\cdot$L$^{-1}$) & 0,020 & 0,020 & 0,040 & 0,010 & 0,020\\
$[\ce{I-}]_0$ (mol$\cdot$L$^{-1}$) & 0,020 & 0,040 & 0,020 & 0,020 & 0,010\\
$t_\mathrm{c}$ (s) & 52 & 26 & 27 & 103 & 105\\\bottomrule
\end{tabular}\end{center}
\begin{questions}
\item Quelle concentration de diiode la réaction (1) a-t-elle produite à la date $t_\mathrm{c}$ ?
\item Montrer que la concentration en ions iodure reste constante jusqu'à $t_\mathrm{c}$ et que celle de l'eau
oxygénée diminue au plus de quelques pourcents. En déduire que la vitesse volumique de formation du diiode par (1)
reste pratiquement constante jusqu'à $t_\mathrm{c}$ et l'exprimer en fonction de $t_\mathrm{c}$.
\item On admet que la vitesse de la réaction (1) vérifie $v=k\,[\ce{H2O2}]\,[\ce{I-}]$. Montrer que les cinq
expériences sont compatibles avec cette loi et déterminer $k$.
\item On propose le mécanisme suivant :
\[\begin{array}{ll}\text{a (lent) : }\ce{H2O2 + I- -> IO- + H2O}\qquad&\text{b (rapide) : }\ce{IO- + H3O+ -> HIO + H2O}\\[1mm]
\text{c (rapide) : }\ce{HIO + I- -> I2 + HO-}\qquad&\text{d (rapide) : }\ce{HO- + H3O+ -> 2 H2O}\end{array}\]
Vérifier qu'il rend compte de l'équation (1) et identifier les intermédiaires réactionnels. On admet que la vitesse
de la transformation est celle de l'acte le plus lent et qu'elle est proportionnelle au produit des concentrations
des entités qui s'y rencontrent : ce mécanisme est-il cohérent avec la loi de la question 3 ?
\item Prévoir $t_\mathrm{c}$ pour $[\ce{H2O2}]_0=[\ce{I-}]_0=\qty{0,030}{mol.L^{-1}}$. Que se passerait-il sans ions
thiosulfate ?
\end{questions}
\end{blocEdition}

\editionPage{defi-ex-6}{D06 · Le générateur de technétium}{Recherche et rédaction}{Très exigeant · 50 à 60~min · Chapitre 6\hfill\hyperref[defi-cor-6-1]{Corrigé p.~\pageref*{defi-cor-6-1}}}
\begin{blocEdition}{defi-ex-6}
Les services de médecine nucléaire produisent eux-mêmes leur technétium 99m grâce à un générateur : du molybdène 99,
fixé sur une colonne, se désintègre en technétium 99m, que l'on extrait par lavage de la colonne (élution).
\begin{center}\small\renewcommand{\arraystretch}{1.2}
\begin{tabular}{@{}lccc@{}}\toprule
Noyau & Demi-vie & Constante radioactive & Nombre de noyaux\\\midrule
molybdène 99 (\ce{^{99}_{42}Mo}) & 66,0~h & $\lambda_1$ & $N_1(t)$\\
technétium 99m (\ce{^{99\text{m}}_{43}Tc}) & 6,01~h & $\lambda_2$ & $N_2(t)$\\\bottomrule
\end{tabular}\end{center}
On admet que chaque désintégration d'un noyau de molybdène 99 donne un noyau de technétium 99m, et qu'à la date
$t=0$ la colonne vient d'être éluée : $N_2(0)=0$ et $N_1(0)=N_{1,0}$.
\begin{questions}
\item Le molybdène 99 est émetteur $\beta^-$. Écrire l'équation de sa désintégration.
\item Justifier que $\dfrac{\mathrm{d}N_2}{\mathrm{d}t}=\lambda_1N_1-\lambda_2N_2$.
\item Vérifier que la fonction $N_2(t)=\dfrac{\lambda_1N_{1,0}}{\lambda_2-\lambda_1}\left(\mathrm{e}^{-\lambda_1t}-\mathrm{e}^{-\lambda_2t}\right)$
est solution de cette équation et qu'elle respecte la condition initiale.
\item Montrer que l'activité $A_2$ du technétium est maximale à la date
$t_\mathrm{m}=\dfrac{\ln(\lambda_2/\lambda_1)}{\lambda_2-\lambda_1}$ et qu'à cette date les activités $A_1$ et $A_2$ sont égales.
Calculer $t_\mathrm{m}$. Pourquoi élue-t-on le générateur environ une fois par jour ?
\item Exprimer le rapport $A_2/A_1$ en fonction du temps. Vers quelle valeur tend-il lorsque $t$ devient grand
devant $t_\mathrm{m}$ ? Interpréter : on parle d'équilibre transitoire.
\end{questions}
En pratique, une partie des désintégrations du molybdène 99 conduit directement au technétium 99 : la simplification
admise ne change pas les conclusions.
\end{blocEdition}

\editionPage{defi-ex-7}{D07 · Le critère d'évolution vu par les vitesses}{Recherche et rédaction}{Très exigeant · 50 à 60~min · Chapitres 4, 5 et 7\hfill\hyperref[defi-cor-7-1]{Corrigé p.~\pageref*{defi-cor-7-1}}}
\begin{blocEdition}{defi-ex-7}
Une transformation \ce{A + B <=> C + D}, en solution, est modélisée par deux actes élémentaires opposés :
\[\ce{A + B -> C + D}\ \text{de vitesse}\ v_1=k_1[\mathrm{A}][\mathrm{B}]\qquad\text{et}\qquad\ce{C + D -> A + B}\ \text{de vitesse}\ v_{-1}=k_{-1}[\mathrm{C}][\mathrm{D}].\]
On donne $k_1=\qty{0,020}{L.mol^{-1}.min^{-1}}$ et $k_{-1}=\qty{0,0050}{L.mol^{-1}.min^{-1}}$. Le volume du mélange vaut
$V=\qty{1,00}{L}$ et reste constant. Le programme ci-dessous simule l'évolution d'un mélange contenant initialement
0,50~mol de A, 0,50~mol de B, 1,50~mol de C et 1,50~mol de D.
\begin{codecote}[0.72]
\begin{pythoncol}
k1, k2, dt = 0.020, 0.005, 0.1
a, b, c, d = 0.50, 0.50, 1.50, 1.50
for i in range(1500):        # 150 min
    v = k1 * a * b - k2 * c * d
    a, b = a - v * dt, b - v * dt
    c, d = c + v * dt, d + v * dt
print(round(a, 3), round(c, 3), round(c*d/(a*b), 2))
\end{pythoncol}
\cote
\begin{consolecol}
0.666 1.334 4.01
\end{consolecol}
\end{codecote}
\begin{questions}
\item Exprimer la vitesse d'apparition $\dfrac{\mathrm{d}[\mathrm{C}]}{\mathrm{d}t}$ en fonction de $v_1$ et de $v_{-1}$.
\item En utilisant le modèle de l'équilibre dynamique, montrer que la constante d'équilibre vaut $K=\dfrac{k_1}{k_{-1}}$.
Calculer $K$.
\item Montrer que $\dfrac{\mathrm{d}[\mathrm{C}]}{\mathrm{d}t}=k_1[\mathrm{A}][\mathrm{B}]\left(1-\dfrac{Q_\mathrm{r}}{K}\right)$. Retrouver ainsi
le critère d'évolution spontanée.
\item Pour le mélange simulé, calculer $Q_{\mathrm{r},\mathrm{i}}$ et prévoir le sens d'évolution. Déterminer par le calcul la
composition du système à l'équilibre, puis la comparer au résultat affiché par le programme.
\item Un catalyseur multiplie $k_1$ et $k_{-1}$ par un même facteur. Quel est son effet sur la constante d'équilibre ?
sur la durée nécessaire pour atteindre l'équilibre ?
\end{questions}
\end{blocEdition}

\editionPage{defi-ex-8}{D08 · Quand une pile est-elle usée ?}{Recherche et rédaction}{Très exigeant · 50 à 60~min · Chapitres 7 et 8\hfill\hyperref[defi-cor-8-1]{Corrigé p.~\pageref*{defi-cor-8-1}}}
\begin{blocEdition}{defi-ex-8}
Une pile Daniell est formée de deux demi-piles de volume $V=\qty{100}{mL}$, de concentrations initiales
$[\ce{Cu^2+}]_0=[\ce{Zn^2+}]_0=c=\qty{0,10}{mol.L^{-1}}$ ; la lame de zinc est en excès. On admet, ce résultat étant hors
programme, que sa tension à vide vaut à 25~°C
\[U=1{,}10-0{,}0296\,\log Q_\mathrm{r}\qquad(U\text{ en volts}),\]
où $Q_\mathrm{r}$ est le quotient de réaction de la réaction de fonctionnement. On donne $F=\qty{96500}{C.mol^{-1}}$ et
$N_\mathrm{A}=\qty{6,02e23}{mol^{-1}}$. Le programme ci-dessous calcule la tension à vide lorsqu'une fraction $f$ de la
capacité $Q_{\max}$ de la pile a été utilisée.
\begin{codecote}[0.72]
\begin{pythoncol}
from math import log10
F, V, c = 96500, 0.100, 0.10    # C/mol, L, mol/L
Qmax = 2 * F * c * V            # capacité (C)
for f in (0.5, 0.9, 0.99, 0.9999):
    q = f * Qmax                # charge débitée (C)
    cu = c - q / (2 * F * V)
    zn = c + q / (2 * F * V)
    print(f, round(1.10 - 0.0296 * log10(zn / cu), 3))
\end{pythoncol}
\cote
\begin{consolecol}
0.5 1.086
0.9 1.062
0.99 1.032
0.9999 0.973
\end{consolecol}
\end{codecote}
\begin{questions}
\item Écrire l'équation de fonctionnement de la pile et l'expression de son quotient de réaction. Calculer la tension à
vide initiale.
\item Montrer que la tension à vide s'annule lorsque $Q_\mathrm{r}=K$. En déduire une valeur de la constante d'équilibre
et la comparer à la valeur $K=\num{1.9e37}$ utilisée au chapitre~8.
\item Exprimer $[\ce{Cu^2+}]$ et $[\ce{Zn^2+}]$ en fonction de la charge $q$ débitée depuis la mise en service.
Calculer la capacité $Q_{\max}$ de la pile.
\item Interpréter les valeurs affichées par le programme. La mesure de la tension à vide permet-elle de savoir qu'une
pile est presque usée ?
\item Quelle serait la concentration des ions \ce{Cu^2+} lorsque la tension à vide s'annule ? À combien d'ions cela
correspond-il dans 100~mL ? Conclure : qu'est-ce qui fixe, en pratique, la durée de vie d'une pile Daniell ?
\end{questions}
\end{blocEdition}

\editionPage{defi-ex-9}{D09 · Le pH d'un acide très dilué}{Recherche et rédaction}{Très exigeant · 40 à 50~min · Chapitres 1, 7 et 9\hfill\hyperref[defi-cor-9-1]{Corrigé p.~\pageref*{defi-cor-9-1}}}
\begin{blocEdition}{defi-ex-9}
On dilue de l'acide chlorhydrique jusqu'à une concentration apportée $c=\qty{1,0e-8}{mol.L^{-1}}$. La relation
$\pH=-\log(c/\cst)$ donnerait $\pH=8$ : la solution serait basique, ce qui est absurde pour un acide. Pour comprendre,
il faut tenir compte des ions apportés par l'eau elle-même. On donne $K_\mathrm{e}=\num{1.0e-14}$ à 25~°C ; le programme
ci-dessous calcule le pH exact d'une solution d'acide fort de concentration $c$.
\begin{codecote}[0.66]
\begin{pythoncol}
from math import log10, sqrt
Ke = 1.0e-14
for c in (1.0e-2, 1.0e-6, 1.0e-8, 1.0e-10):
    h = (c + sqrt(c**2 + 4*Ke)) / 2
    print(c, round(-log10(h), 2))
\end{pythoncol}
\cote
\begin{consolecol}
0.01 2.0
1e-06 6.0
1e-08 6.98
1e-10 7.0
\end{consolecol}
\end{codecote}
\begin{questions}
\item Expliquer sans calcul pourquoi l'ajout d'un acide, même infime, ne peut pas rendre une solution basique.
\item La solution contient les ions \ce{H3O+}, \ce{HO-} et \ce{Cl-}, avec $[\ce{Cl-}]=c$. En écrivant qu'elle est
électriquement neutre, montrer que $h=[\ce{H3O+}]_\text{éq}$ vérifie $h^2-c\,h-K_\mathrm{e}=0$ (concentrations en
\unit{mol.L^{-1}}).
\item Résoudre cette équation et retrouver la ligne du programme qui calcule $h$. Vérifier la valeur du pH pour
$c=\qty{1,0e-8}{mol.L^{-1}}$.
\item Interpréter les quatre lignes affichées : pour quelles concentrations la relation usuelle $\pH=-\log(c/\cst)$
est-elle fiable ? Que devient le pH lorsque $c$ tend vers 0 ?
\item Sans refaire tous les calculs, décrire ce que donnerait l'étude analogue pour une base forte très diluée.
\end{questions}
\end{blocEdition}

\editionPage{defi-ex-10}{D10 · Pourquoi au moins 1,23 volt ?}{Recherche et rédaction}{Très exigeant · 40 à 50~min · Chapitres 8 et 10\hfill\hyperref[defi-cor-10-1]{Corrigé p.~\pageref*{defi-cor-10-1}}}
\begin{blocEdition}{defi-ex-10}
Aussi longtemps que la tension d'une cellule d'électrolyse de l'eau reste inférieure à 1,23~V, il ne se passe rien ;
au-delà, le dihydrogène et le dioxygène apparaissent. Ce seuil n'est pas un hasard : 1,23~V est la tension à vide de la
pile à combustible dihydrogène-dioxygène, dont l'électrolyse de l'eau est la transformation inverse. On donne
$F=\qty{96500}{C.mol^{-1}}$ et l'énergie restituée par la combustion d'une mole de dihydrogène : $\qty{242}{kJ}$.
\begin{codecote}[0.6]
\begin{pythoncol}
F, U0 = 96500, 1.23
for U in (1.3, 1.6, 2.0, 2.4):
    w = 2 * F * U / 1000   # kJ par mole
    print(U, round(w), round(U0 / U, 2))
\end{pythoncol}
\cote
\begin{consolecol}
1.3 251 0.95
1.6 309 0.77
2.0 386 0.61
2.4 463 0.51
\end{consolecol}
\end{codecote}
\begin{questions}
\item Dans la cellule, les deux électrodes couvertes de gaz constituent une pile dihydrogène-dioxygène de tension à
vide $U_0=\qty{1,23}{V}$. Expliquer pourquoi aucune électrolyse durable n'est possible si le générateur impose une
tension inférieure à $U_0$.
\item Montrer que l'énergie électrique minimale nécessaire pour produire une mole de dihydrogène vaut
$W_0=2F\,U_0\approx\qty{237}{kJ}$.
\item En déduire l'énergie électrique minimale par kilogramme de dihydrogène, en kilowattheures
($\qty{1}{kW.h}=\qty{3,6e6}{J}$).
\item Interpréter les trois colonnes affichées par le programme. Que représente le quotient $U_0/U$ et comment
évolue-t-il ?
\item Comparer $W_0$ à l'énergie restituée par la combustion d'une mole de dihydrogène. Qu'est-ce que ce résultat,
surprenant, impose à une cellule qui fonctionnerait exactement à 1,23~V ? À quelle tension la cellule n'échange-t-elle
aucune chaleur avec l'extérieur ?
\end{questions}
\end{blocEdition}

\editionPage{defi-ex-11}{D11 · Combien d'isomères ?}{Recherche et rédaction}{Avancé · 40 à 50~min · Chapitre 11\hfill\hyperref[defi-cor-11-1]{Corrigé p.~\pageref*{defi-cor-11-1}}}
\begin{blocEdition}{defi-ex-11}
Les alcanes ont pour formule brute $\ce{C_{$n$}H_{$2n+2$}}$. Pour $n\leq3$, un seul enchaînement est possible ; pour
$n=4$, il en existe deux (butane et 2-méthylpropane). Le but du défi est de dénombrer, sans en oublier ni en compter
deux fois, tous les squelettes possibles pour $n=5$ puis $n=6$ — puis de prendre la mesure de ce que devient ce
dénombrement pour les grandes molécules : on compte 75 isomères de constitution pour \ce{C10H22} et plus de
\num{366000} pour \ce{C20H42}.
\begin{questions}
\item Rappeler ce que sont des isomères de constitution. Justifier qu'un alcane à $n$ atomes de carbone, quelle que
soit la forme de son squelette non cyclique, possède exactement $n-1$ liaisons C--C.
\item Dessiner tous les squelettes possibles pour \ce{C5H12}, en expliquant la méthode suivie et pourquoi aucun autre
n'existe. Nommer chaque isomère.
\item Faire de même pour \ce{C6H14}.
\item Au cours de la recherche, on obtient parfois deux dessins d'apparence différente qui représentent le même
squelette. Donner un exemple rencontré pour \ce{C6H14} et expliquer comment la nomenclature permet de démasquer ces
doublons.
\item À l'aide des valeurs données pour \ce{C10H22} et \ce{C20H42}, expliquer pourquoi la chimie organique ne peut se
passer ni des formules topologiques, ni d'une nomenclature systématique.
\end{questions}
\end{blocEdition}

\editionPage{defi-ex-12}{D12 · Jusqu'où pousser un rendement ?}{Recherche et rédaction}{Très exigeant · 50 à 60~min · Chapitres 7 et 12\hfill\hyperref[defi-cor-12-1]{Corrigé p.~\pageref*{defi-cor-12-1}}}
\begin{blocEdition}{defi-ex-12}
L'estérification \ce{CH3COOH + C2H5OH <=> CH3COOC2H5 + H2O} a pour constante d'équilibre $K=4{,}0$ et, en mélange
équimolaire, plafonne à 67~\% d'avancement. Pour gagner les points manquants, le chimiste introduit l'alcool en
excès : on note $r=n_{\text{alcool}}/n_{\text{acide}}$ le rapport initial, et $x$ le taux d'avancement final de
l'acide. Le programme ci-dessous cherche l'excès nécessaire pour atteindre 95~\%.
\begin{codecote}[0.66]
\begin{pythoncol}
from math import sqrt
def taux(r):
    d = sqrt(r*r - r + 1)
    return (2/3) * (1 + r - d)
r = 1.0
while taux(r) < 0.95:
    r = r + 0.5
print(r, round(taux(r), 3))
\end{pythoncol}
\cote
\begin{consolecol}
5.5 0.95
\end{consolecol}
\end{codecote}
\begin{questions}
\item À l'aide du chapitre~7, rappeler pourquoi, en mélange équimolaire et quelle que soit la durée de chauffage, le
taux d'avancement final vaut $2/3$.
\item Pour 1~mol d'acide et $r$~mol d'alcool, dresser le tableau d'avancement et montrer que le taux d'avancement
final $x$ vérifie $3x^2-4(1+r)\,x+4r=0$.
\item En déduire que $x=\dfrac{2}{3}\left(1+r-\sqrt{r^2-r+1}\right)$, puis vérifier les valeurs attendues pour $r=1$
et $r=2$ (exercice~10 du chapitre~7).
\item Interpréter le fonctionnement et la sortie du programme. Quel excès faut-il pour atteindre 95~\% ? Commenter
son coût.
\item Comparer cette stratégie à l'élimination de l'eau au fur et à mesure. Pourquoi un catalyseur, lui, ne
modifie-t-il jamais la valeur du palier ?
\end{questions}
\end{blocEdition}

\editionPage{defi-ex-13}{D13 · Le caillou coincé dans un pneu}{Recherche et rédaction}{Très exigeant · 45 à 60~min · Chapitre 13\hfill\hyperref[defi-cor-13-1]{Corrigé p.~\pageref*{defi-cor-13-1}}}
\begin{blocEdition}{defi-ex-13}
Un vélo roule en ligne droite à la vitesse constante $V=\qty{8,0}{m.s^{-1}}$ ; ses roues, de rayon $R=\qty{0,35}{m}$, roulent
sans glisser. Un petit caillou est coincé dans la bande de roulement d'un pneu. Dans le référentiel terrestre, avec
$(Ox)$ horizontal dans le sens du mouvement et $(Oy)$ vertical vers le haut, ses coordonnées sont, en posant
$\omega=V/R$ :
\[x(t)=R\left(\omega t-\sin\omega t\right)\qquad\text{et}\qquad y(t)=R\left(1-\cos\omega t\right).\]
Sa trajectoire est une \textbf{cycloïde}. Le programme ci-dessous calcule la norme de l'accélération du caillou à trois
dates.
\begin{codecote}[0.62]
\begin{pythoncol}
from math import sin, cos, sqrt
R, V = 0.35, 8.0
w = V / R
for t in (0.00, 0.05, 0.10):
    ax = R * w**2 * sin(w * t)
    ay = R * w**2 * cos(w * t)
    print(t, round(sqrt(ax**2 + ay**2)))
\end{pythoncol}
\cote
\begin{consolecol}
0.0 183
0.05 183
0.1 183
\end{consolecol}
\end{codecote}
\begin{questions}
\item Le centre C de la roue a pour coordonnées $x_\mathrm{C}=R\,\omega t$ et $y_\mathrm{C}=R$. Décrire son mouvement.
\item Établir les coordonnées du vecteur vitesse du caillou. Calculer sa vitesse lorsqu'il touche le sol, puis
lorsqu'il est au sommet de la roue. Commenter.
\item Établir les coordonnées de son vecteur accélération et justifier les lignes correspondantes du programme.
\item Montrer que la norme de l'accélération est constante, et retrouver la valeur affichée.
\item \textit{Défi.} Montrer que le vecteur accélération est toujours dirigé vers le centre C de la roue. Le caillou
a-t-il un mouvement circulaire uniforme dans le référentiel terrestre ? dans un référentiel lié au vélo ?
\end{questions}
\end{blocEdition}

\editionPage{defi-ex-14}{D14 · De quelle hauteur lancer le wagon ?}{Recherche et rédaction}{Très exigeant · 45 à 60~min · Chapitres 13 et 14\hfill\hyperref[defi-cor-14-1]{Corrigé p.~\pageref*{defi-cor-14-1}}}
\begin{blocEdition}{defi-ex-14}
Un wagon de masse $m$, assimilé à un point, est lâché sans vitesse du haut d'une rampe, à la hauteur $h$ au-dessus du
point bas d'un looping \textbf{circulaire} de rayon $R=\qty{6,0}{m}$. On néglige tous les frottements : l'énergie
mécanique du wagon se conserve (classe de première). Sur le looping, le rail exerce sur le wagon une force $\vec{N}$
perpendiculaire au rail, dirigée vers le centre du cercle ; le rail ne peut que pousser le wagon, si bien que
$N\geq0$. Le programme ci-dessous compare $N$ au poids, au sommet et au point bas, pour trois hauteurs de départ.
\begin{codecote}[0.66]
\begin{pythoncol}
g, R = 9.81, 6.0
for h in (12.0, 15.0, 18.0):
    v2_haut = 2 * g * (h - 2 * R)
    n_haut = v2_haut / (g * R) - 1   # N/(mg)
    n_bas = 1 + 2 * h / R            # N/(mg)
    print(h, round(n_haut, 2), round(n_bas, 2))
\end{pythoncol}
\cote
\begin{consolecol}
12.0 -1.0 5.0
15.0 0.0 6.0
18.0 1.0 7.0
\end{consolecol}
\end{codecote}
\begin{questions}
\item En utilisant la conservation de l'énergie mécanique, exprimer le carré de la vitesse du wagon au sommet du
looping en fonction de $g$, $h$ et $R$.
\item Appliquer la deuxième loi de Newton au wagon au sommet du looping, dans le repère de Frenet. En déduire
l'expression de $N/(m\,g)$ au sommet et retrouver la ligne correspondante du programme.
\item Quelle condition $N$ doit-il vérifier pour que le wagon reste sur le rail ? En déduire la hauteur minimale de
départ, puis interpréter la première ligne affichée.
\item Établir de même l'expression de $N/(m\,g)$ au point bas. Que vaut-elle pour la hauteur minimale ?
\item À l'aide de l'exercice~14 du chapitre~13, expliquer pourquoi les loopings réels ne sont pas circulaires.
\end{questions}
\end{blocEdition}

\editionPage{defi-ex-15}{D15 · Lancer à 45° ?}{Recherche et rédaction}{Avancé · 40 à 50~min · Chapitre 15\hfill\hyperref[defi-cor-15-1]{Corrigé p.~\pageref*{defi-cor-15-1}}}
\begin{blocEdition}{defi-ex-15}
« Pour lancer le plus loin possible, il faut lancer à 45° » : la règle est souvent citée. Elle est exacte pour un
projectile lancé depuis le sol, sans frottement. Un lanceur de poids lâche pourtant l'engin à environ \qty{2,10}{m} de
hauteur. On reprend les données de l'exercice~6 : $v_0=\qty{13,5}{m.s^{-1}}$ et $g=\qty{9,81}{m.s^{-2}}$. Le programme
ci-dessous calcule la longueur du jet pour tous les angles entiers de $30^\circ$ à $50^\circ$ et retient le meilleur.
\begin{codecote}[0.66]
\begin{pythoncol}
from math import sin, cos, sqrt, radians
v0, g, h = 13.5, 9.81, 2.10
meilleur, angle = 0, 0
for a in range(30, 51):
    vx = v0 * cos(radians(a))
    vz = v0 * sin(radians(a))
    t = (vz + sqrt(vz**2 + 2*g*h)) / g
    if vx * t > meilleur:
        meilleur, angle = vx * t, a
print(angle, round(meilleur, 2))
\end{pythoncol}
\cote
\begin{consolecol}
42 20.57
\end{consolecol}
\end{codecote}
\begin{questions}
\item Pour un lancer depuis le sol ($h=0$), montrer que la portée vaut $D=\dfrac{v_0^2\sin(2\alpha)}{g}$. En déduire l'angle qui
la rend maximale et la valeur correspondante.
\item Pour un lancer depuis la hauteur $h$, justifier l'expression de la durée du vol utilisée par le programme.
\item Interpréter la sortie du programme. Pourquoi l'angle optimal est-il inférieur à $45^\circ$ lorsque $h>0$ ?
\item Comparer les longueurs obtenues à $42^\circ$ et à $35^\circ$. Les meilleurs lanceurs lâchent le poids vers $37^\circ$ :
sachant qu'ils communiquent une vitesse plus grande quand l'angle est plus faible, commenter.
\end{questions}
\end{blocEdition}

\editionPage{defi-ex-16}{D16 · Peser le trou noir de la Galaxie}{Recherche et rédaction}{Très exigeant · 50 à 60~min · Chapitre 16\hfill\hyperref[defi-cor-16-1]{Corrigé p.~\pageref*{defi-cor-16-1}}}
\begin{blocEdition}{defi-ex-16}
Au centre de la Voie lactée, des étoiles tournent à grande vitesse autour d'un objet invisible, Sagittarius~A*. Leur
suivi pendant trente ans a valu le prix Nobel de physique 2020 à Reinhard Genzel et Andrea Ghez. Les orbites de trois de
ces étoiles sont des ellipses de même foyer ; leurs demi-grands axes $a$ et leurs périodes $T$ valent, en arrondissant :
S2 : \qty{1031}{ua} et 16,05~ans ; S38 : \qty{1169}{ua} et 19,2~ans ; S55 : \qty{889}{ua} et 12,8~ans.
\begin{codecote}[0.62]
\begin{pythoncol}
etoiles = {"S2": (1031, 16.05),
           "S38": (1169, 19.2),
           "S55": (889, 12.8)}
for nom, (a, T) in etoiles.items():
    M = a**3 / T**2          # en masses solaires
    print(nom, round(M / 1e6, 2))
\end{pythoncol}
\cote
\begin{consolecol}
S2 4.25
S38 4.33
S55 4.29
\end{consolecol}
\end{codecote}
\par\smallskip
Données : $G=\qty{6,67e-11}{N.m^2.kg^{-2}}$ ; masse du Soleil $M_\odot=\qty{1,99e30}{kg}$ ; $\qty{1}{ua}=\qty{1,496e11}{m}$ ;
$c=\qty{3,00e8}{m.s^{-1}}$ ; excentricité de l'orbite de S2 : $e=0{,}886$ (distances au foyer extrêmes : $a(1-e)$ et
$a(1+e)$).
\begin{questions}
\item Pour la Terre autour du Soleil, $T=\qty{1}{an}$ et $a=\qty{1}{ua}$. En écrivant la troisième loi de Kepler pour la Terre et
pour une étoile tournant autour d'un astre de masse $M$, montrer que $\dfrac{M}{M_\odot}=\dfrac{a^3}{T^2}$, $a$ étant en ua et $T$ en
années. Justifier la ligne de calcul du programme.
\item Interpréter les valeurs affichées. Que peut-on conclure sur l'objet central ?
\item Calculer sa masse en kilogrammes. On admet que le rayon d'un trou noir de masse $M$ vaut $R_\mathrm{S}=\dfrac{2GM}{c^2}$ :
le calculer, en ua.
\item Calculer la distance minimale entre S2 et l'objet central. Combien de fois est-elle plus grande que $R_\mathrm{S}$ ?
Pourquoi cette observation renforce-t-elle l'hypothèse d'un trou noir ?
\item En appliquant la loi des aires aux deux extrémités du grand axe, calculer le rapport des vitesses de S2 au plus près
et au plus loin. Au plus près, S2 file à \qty{7650}{km.s^{-1}} : quelle est sa vitesse au plus loin ?
\end{questions}
\end{blocEdition}

\editionPage{defi-ex-17}{D17 · La clepsydre}{Recherche et rédaction}{Avancé · 45 à 60~min · Chapitre 17\hfill\hyperref[defi-cor-17-1]{Corrigé p.~\pageref*{defi-cor-17-1}}}
\begin{blocEdition}{defi-ex-17}
Dans l'Égypte ancienne, on mesurait le temps, la nuit, avec des clepsydres : des vases percés d'un petit orifice près du
fond, dont le niveau d'eau baisse pendant la vidange. La plus ancienne connue, trouvée à Karnak, date d'environ 1400
av.~J.-C. Pour que des graduations régulièrement espacées marquent des durées égales, le niveau doit baisser à vitesse
constante. On étudie un vase de révolution de hauteur $H=\qty{30}{cm}$, de rayon $R=\qty{10}{cm}$ au sommet, qui doit se
vider en 12~heures. On note $h$ la hauteur d'eau, $r$ le rayon du vase à cette hauteur et $s$ la section de l'orifice ;
l'eau sort à la vitesse $v=\sqrt{2\,g\,h}$ (formule de Torricelli), avec $g=\qty{9,81}{m.s^{-2}}$. Le programme simule la
vidange d'un vase de profil $r=R\,(h/H)^p$, avec le même orifice, pour un cylindre ($p=0$) et pour la clepsydre
($p=1/4$) : il affiche les hauteurs d'eau, en centimètres, toutes les trois heures.
\begin{codecote}[0.66]
\begin{pythoncol}
from math import pi, sqrt
g, H, R = 9.81, 0.30, 0.10
s = 9.0e-8                    # orifice (m²)
def niveau(p, duree):
    h, t, dt = H, 0, 1
    while t < duree and h > 0:
        r = R * (h / H)**p
        v = sqrt(2 * g * h)
        h = h - s*v / (pi*r**2) * dt
        t = t + dt
    return round(100 * max(h, 0.0), 1)
for n in range(0, 13, 3):
    T = 3600 * n
    print(n, niveau(0, T), niveau(0.25, T))
\end{pythoncol}
\cote
\begin{consolecol}
0 30.0 30.0
3 23.0 22.5
6 16.9 15.0
9 11.7 7.5
12 7.5 0.0
\end{consolecol}
\end{codecote}
\par\smallskip
\begin{questions}
\item Pendant la durée $\mathrm{d}t$, le niveau baisse de $|\mathrm{d}h|$. En écrivant que le volume perdu par le vase est sorti par
l'orifice, montrer que le niveau descend à la vitesse $u=-\dfrac{\mathrm{d}h}{\mathrm{d}t}=\dfrac{s\,\sqrt{2\,g\,h}}{\pi\,r(h)^2}$.
\item Pour un cylindre de rayon $R$, vérifier que $\sqrt{h}=\sqrt{H}-\dfrac{s}{\pi R^2}\sqrt{\dfrac{g}{2}}\;t$ convient. En déduire la durée
de vidange. Pourquoi des graduations régulières seraient-elles fausses ?
\item Quel profil donne une vitesse de descente $u$ constante ? Montrer que $h=H\,(r/R)^4$. Calculer $u$, puis $s$ et le
diamètre de l'orifice pour une vidange en 12~heures.
\item Interpréter la sortie du programme.
\item Montrer que le cylindre se vide en exactement deux fois plus de temps que la clepsydre. Les clepsydres égyptiennes
avaient la forme d'un tronc de cône évasé vers le haut : commenter.
\end{questions}
\end{blocEdition}

\editionPage{defi-ex-18}{D18 · Faire le vide}{Recherche et rédaction}{Avancé · 40 à 50~min · Chapitre 18\hfill\hyperref[defi-cor-18-1]{Corrigé p.~\pageref*{defi-cor-18-1}}}
\begin{blocEdition}{defi-ex-18}
Pour faire le vide dans un récipient de volume $V=\qty{1,0}{L}$, initialement rempli d'air sous $p_0=\qty{1013}{hPa}$, on
utilise une pompe à piston dont le cylindre a un volume $v=\qty{50}{mL}$. À chaque coup de pompe, le piston recule : l'air
du récipient se détend dans le récipient et le cylindre réunis. Une soupape se ferme, puis le piston refoule dans
l'atmosphère l'air contenu dans le cylindre. La température reste constante et l'air se comporte comme un gaz parfait.
Le programme simule les coups de pompe successifs jusqu'à ce que la pression descende sous \qty{1}{hPa} ; \code{vm} est le
volume mort, partie du cylindre que le piston ne peut pas vider.
\begin{codecote}[0.62]
\begin{pythoncol}
p0, V, v, vm = 1013, 1.0, 0.050, 0.0  # hPa, L
p, k = p0, 0
while p > 1 and k < 1000:
    p = (p*V + p0*vm) / (V + v + vm)
    k = k + 1
print(k, round(p, 3))
\end{pythoncol}
\cote
\begin{consolecol}
142 0.992
\end{consolecol}
\end{codecote}
\par\smallskip
\begin{questions}
\item Sans volume mort, montrer qu'après le premier coup de pompe la pression vaut $p_1=p_0\,\dfrac{V}{V+v}$. La calculer.
\item Montrer que la suite des pressions $(p_k)$ est géométrique et exprimer $p_k$ en fonction de $k$.
\item Retrouver par le calcul le nombre de coups affiché par le programme.
\item Pour gonfler un pneu de vélo de \qty{1,5}{L}, initialement sous \qty{1,0}{bar}, chaque coup d'une pompe de \qty{0,10}{L}
y fait entrer l'air atmosphérique contenu dans la pompe. Comment évolue la pression du pneu au fil des coups ?
Comparer avec la mise sous vide.
\item Une pompe réelle a un volume mort $v_\mathrm{m}=\qty{0,50}{mL}$ : après chaque refoulement, il y reste de l'air sous
$p_0$. Justifier la ligne de calcul du programme, puis montrer que la pression ne peut pas descendre sous
$p_\mathrm{lim}=p_0\,\dfrac{v_\mathrm{m}}{v+v_\mathrm{m}}$. Calculer $p_\mathrm{lim}$ et prévoir la sortie du programme avec
\code{vm = 0.0005}.
\end{questions}
\end{blocEdition}

\editionPage{defi-ex-19}{D19 · L'expérience de Joule}{Recherche et rédaction}{Avancé · 40 à 50~min · Chapitre 19\hfill\hyperref[defi-cor-19-1]{Corrigé p.~\pageref*{defi-cor-19-1}}}
\begin{blocEdition}{defi-ex-19}
En 1850, James Prescott Joule publie une expérience restée célèbre : la chute de masses fait tourner des palettes qui
brassent l'eau d'un récipient en cuivre, dont la température s'élève légèrement. Il montre ainsi qu'un travail produit
le même effet qu'un transfert thermique. On en étudie une version simplifiée, aux données arrondies : deux masses de
\qty{13,0}{kg} chacune descendent de \qty{1,60}{m}, puis sont remontées sans entraîner les palettes ; l'opération est
répétée 20 fois. Les masses arrivent en bas à la vitesse de \qty{0,060}{m.s^{-1}}. Le récipient contient \qty{6,00}{kg}
d'eau ; le récipient et les palettes ont une capacité thermique de \qty{400}{J.K^{-1}}.
\schema{19-joule}
Données : $g=\qty{9,81}{N.kg^{-1}}$ ; $c_\mathrm{eau}=\qty{4180}{J.kg^{-1}.K^{-1}}$.
\begin{questions}
\item On choisit le système \{eau, récipient, palettes\}. Quel transfert d'énergie reçoit-il pendant une descente ?
Pourquoi sa température augmente-t-elle, alors qu'aucun corps chaud ne le touche ?
\item On admet que le travail reçu par le système est égal à l'énergie mécanique perdue par les masses. Le calculer
pour les 20 descentes et montrer que le terme dû à la vitesse finale des masses est négligeable.
\item En négligeant les échanges avec l'air, prévoir l'élévation de température du système.
\item Joule mesure une élévation de \qty{0,313}{K}. À son époque, les transferts thermiques se mesuraient en calories :
\qty{1}{cal} élève de \qty{1}{\degreeCelsius} la température de \qty{1}{g} d'eau. Montrer que le récipient et les palettes
équivalent à \qty{95,7}{g} d'eau, puis calculer combien de joules correspondent, d'après ces mesures, à une calorie.
Comparer à la valeur actuelle, \qty{4,18}{J}.
\item Proposer des causes à l'écart. Pourquoi Joule répétait-il vingt fois la descente ? En quoi cette expérience
fonde-t-elle le premier principe ?
\end{questions}
\end{blocEdition}

\editionPage{defi-ex-20}{D20 · Café : le lait avant ou après ?}{Recherche et rédaction}{Avancé · 40 à 50~min · Chapitre 20\hfill\hyperref[defi-cor-20-1]{Corrigé p.~\pageref*{defi-cor-20-1}}}
\begin{blocEdition}{defi-ex-20}
Vous vous servez \qty{200}{g} de café à \qty{90}{\degreeCelsius}, dans une pièce à \qty{20}{\degreeCelsius}, et vous comptez le boire
dans dix minutes avec \qty{50}{g} de lait sortant du réfrigérateur, à \qty{5}{\degreeCelsius}. Pour obtenir le café le plus frais,
faut-il verser le lait tout de suite (option A) ou juste avant de boire (option B) ? On admet que la boisson refroidit
selon la loi de Newton, avec un temps caractéristique $\tau=\qty{25}{min}$, le même avec ou sans lait ; le lait et le café
ont la même capacité thermique massique ; le mélange est instantané et n'échange pas d'énergie avec l'extérieur.
\begin{codecote}[0.66]
\begin{pythoncol}
from math import exp
Te, Tc, Tl, tau = 20, 90, 5, 1500   # °C, s
mc, ml, t = 200, 50, 600            # g, s
def refroidir(T0, t):
    return Te + (T0 - Te) * exp(-t / tau)
def melanger(Tc, Tl):
    return (mc * Tc + ml * Tl) / (mc + ml)
TA = refroidir(melanger(Tc, Tl), t)
TB = melanger(refroidir(Tc, t), Tl)
print(round(TA, 2), round(TB, 2))
\end{pythoncol}
\cote
\begin{consolecol}
55.53 54.54
\end{consolecol}
\end{codecote}
\par\smallskip
\begin{questions}
\item Par un bilan d'énergie, montrer que deux liquides de masses $m_\mathrm{c}$ et $m_\mathrm{l}$, de températures $T_\mathrm{c}$ et
$T_\mathrm{l}$, atteignent la température $\dfrac{m_\mathrm{c}T_\mathrm{c}+m_\mathrm{l}T_\mathrm{l}}{m_\mathrm{c}+m_\mathrm{l}}$. Quelle fonction du programme
utilise ce résultat ?
\item Justifier la fonction \code{refroidir}, puis calculer à la main la température $T_\mathrm{A}$ du café de l'option A.
\item Calculer la température $T_\mathrm{B}$ de l'option B. Quelle option donne le café le plus frais ? Proposer une
explication physique.
\item Montrer que $T_\mathrm{A}-T_\mathrm{B}=\dfrac{m_\mathrm{l}}{m_\mathrm{c}+m_\mathrm{l}}\,(T_\mathrm{e}-T_\mathrm{l})\left(1-\mathrm{e}^{-t/\tau}\right)$.
Que se passe-t-il si le lait est à la température de la pièce ? Quel écart maximal peut-on obtenir ?
\item Le lait augmente la capacité thermique du contenu de la tasse. Comment cela modifie-t-il le temps caractéristique ?
La conclusion en est-elle affaiblie ou renforcée ?
\end{questions}
\end{blocEdition}

\editionPage{defi-ex-21}{D21 · Pourquoi entend-on les basses du voisin ?}{Recherche et rédaction}{Avancé · 40 à 50~min · Chapitre 21\hfill\hyperref[defi-cor-21-1]{Corrigé p.~\pageref*{defi-cor-21-1}}}
\begin{blocEdition}{defi-ex-21}
Lorsqu'un son frappe une paroi, une partie de son énergie est réfléchie, une partie absorbée, une partie transmise.
L'\textbf{indice d'affaiblissement} $R$ d'une paroi est l'atténuation, en décibels, entre le niveau sonore du côté de la
source et le niveau transmis de l'autre côté. Pour une paroi simple, la « loi de masse » en donne une estimation :
$R=20\log(m_\mathrm{s}\,f)-47$, où $m_\mathrm{s}$ est la masse surfacique de la paroi (en \unit{kg.m^{-2}}) et $f$ la fréquence
du son (en Hz). Le programme calcule $R$ pour un mur de béton de \qty{10}{cm}, de masse surfacique \qty{230}{kg.m^{-2}}.
\begin{codecote}[0.66]
\begin{pythoncol}
from math import log10
ms = 230                     # kg/m², béton de 10 cm
for f in [125, 250, 500, 1000, 2000, 4000]:   # Hz
    R = 20 * log10(ms * f) - 47   # loi de masse
    print(f, round(R, 1))
\end{pythoncol}
\cote
\begin{consolecol}
125 42.2
250 48.2
500 54.2
1000 60.2
2000 66.3
4000 72.3
\end{consolecol}
\end{codecote}
\par\smallskip
Données : cloison faite de deux plaques de plâtre, $m_\mathrm{s}=\qty{20}{kg.m^{-2}}$ ; porte en bois, $m_\mathrm{s}=\qty{15}{kg.m^{-2}}$.
\begin{questions}
\item Montrer que $R$ augmente de \qty{6}{dB} chaque fois que la fréquence double. Le vérifier sur la sortie du programme.
\item Une conversation produit \qty{60}{dB} à \qty{500}{Hz} contre une paroi. Quel niveau entend-on de l'autre côté du mur de
béton ? de la cloison en plaques de plâtre ? Conclure.
\item Le voisin écoute de la musique : \qty{90}{dB} à \qty{63}{Hz} pour les basses, \qty{80}{dB} à \qty{4000}{Hz} pour les aigus.
Calculer les niveaux transmis par le mur de béton et expliquer le titre du défi.
\item Par combien faudrait-il multiplier la masse du mur pour gagner \qty{20}{dB} à toutes les fréquences ? Est-ce réaliste ?
\item Le mur de \qty{10}{m^2} comporte une porte en bois de \qty{2}{m^2}. Une atténuation $R$ laisse passer la fraction
$\tau=10^{-R/10}$ de l'intensité, et les fractions transmises par chaque partie s'ajoutent en proportion des surfaces.
Calculer l'indice d'affaiblissement de l'ensemble à \qty{500}{Hz} et commenter.
\end{questions}
\end{blocEdition}

\editionPage{defi-ex-22}{D22 · Des franges dans la tache de diffraction}{Recherche et rédaction}{Avancé · 45 à 60~min · Chapitre 22\hfill\hyperref[defi-cor-22-1]{Corrigé p.~\pageref*{defi-cor-22-1}}}
\begin{blocEdition}{defi-ex-22}
Dans une expérience de fentes d'Young réelle, chaque fente a une largeur $a$ non négligeable et diffracte la lumière :
les franges d'interférences n'apparaissent que là où parvient la lumière diffractée. On admet qu'au point d'abscisse $x$
de l'écran, l'intensité lumineuse relative est le produit d'un terme de diffraction $\left(\dfrac{\sin u}{u}\right)^2$, avec
$u=\dfrac{\pi\,a\,x}{\lambda D}$, et d'un terme d'interférences $\cos^2\left(\dfrac{\pi\,b\,x}{\lambda D}\right)$. On utilise un laser de
longueur d'onde $\lambda=\qty{633}{nm}$, des fentes de largeur $a=\qty{0,10}{mm}$ distantes de $b=\qty{0,50}{mm}$ et un écran à
$D=\qty{2,0}{m}$. Le programme calcule l'intensité relative au centre des franges brillantes successives.
\begin{codecote}[0.7]
\begin{pythoncol}
from math import sin, cos, pi
lam, D, a, b = 633e-9, 2.0, 0.10e-3, 0.50e-3
def intensite(x):
    u = pi * a * x / (lam * D)
    env = 1 if u == 0 else (sin(u) / u)**2
    return env * cos(pi * b * x / (lam * D))**2
i = lam * D / b                  # interfrange
for k in range(6):
    x = k * i
    I = round(intensite(x), 3)
    print(k, round(1000 * x, 2), I)
\end{pythoncol}
\cote
\begin{consolecol}
0 0.0 1.0
1 2.53 0.875
2 5.06 0.573
3 7.6 0.255
4 10.13 0.055
5 12.66 0.0
\end{consolecol}
\end{codecote}
\par\smallskip
\begin{questions}
\item Avec une seule fente, calculer la largeur $L$ de la tache centrale de diffraction.
\item Calculer l'interfrange $i$ et montrer que $\dfrac{L}{i}=\dfrac{2b}{a}$.
\item Où se trouve le premier zéro de la figure de diffraction ? Montrer que la frange brillante d'ordre 5 tombe sur ce
zéro. Combien de franges brillantes la tache centrale contient-elle ?
\item Interpréter la sortie du programme. Pourquoi l'intensité des franges décroît-elle en s'éloignant du centre ?
\item Avec des fentes deux fois plus fines, la distance $b$ restant la même, combien de franges brillantes compterait-on
dans la tache centrale ? Quel inconvénient cela présente-t-il ?
\end{questions}
\end{blocEdition}

\editionPage{defi-ex-23}{D23 · Le passage d'une voiture}{Recherche et rédaction}{Avancé · 45 à 60~min · Chapitre 23\hfill\hyperref[defi-cor-23-1]{Corrigé p.~\pageref*{defi-cor-23-1}}}
\begin{blocEdition}{defi-ex-23}
Un auditeur se tient à la distance $b=\qty{10}{m}$ d'une route rectiligne. Une voiture passe à $v_\mathrm{E}=\qty{25}{m.s^{-1}}$ en
klaxonnant à $f_\mathrm{E}=\qty{900}{Hz}$ ; le son se propage à $v=\qty{340}{m.s^{-1}}$. Lorsque la voiture n'est pas dans l'axe de
l'auditeur, seule compte la composante de sa vitesse dans la direction de l'auditeur : $f_\mathrm{R}=f_\mathrm{E}\,\dfrac{v}{v-v_\mathrm{E}\cos\theta}$,
où $\theta$ est l'angle entre la vitesse de la voiture et la direction voiture-auditeur au moment de l'émission. On
repère la voiture par son abscisse $x=v_\mathrm{E}\,t$ le long de la route, l'origine étant le point le plus proche de
l'auditeur. Le programme affiche, pour plusieurs dates d'émission $t$, la date de réception et la fréquence reçue.
\begin{codecote}[0.66]
\begin{pythoncol}
from math import sqrt
v, fE, vE, b = 340.0, 900.0, 25.0, 10.0
for t in [-4, -2, -1, 0, 1, 2, 4]:
    x = vE * t               # abscisse (m)
    d = sqrt(x**2 + b**2)    # distance (m)
    cos_theta = -x / d
    fR = fE * v / (v - vE * cos_theta)
    print(t, round(t + d/v, 2), round(fR, 1))
\end{pythoncol}
\cote
\begin{consolecol}
-4 -3.7 971.0
-2 -1.85 969.9
-1 -0.92 965.9
0 0.03 900.0
1 1.08 842.5
2 2.15 839.5
4 4.3 838.6
\end{consolecol}
\end{codecote}
\par\smallskip
\begin{questions}
\item Que devient $\cos\theta$ très longtemps avant et très longtemps après le passage ? Montrer que $f_\mathrm{R}$ tend alors
vers les valeurs données par les relations à une dimension, et les calculer.
\item Quelle fréquence l'auditeur reçoit-il pour le son émis au moment où la voiture passe au plus près de lui ? À quelle
date le reçoit-il ?
\item Pourquoi les dates de réception diffèrent-elles des dates d'émission ?
\item Le décalage n'est pas symétrique : comparer $f_\mathrm{app}-f_\mathrm{E}$ et $f_\mathrm{E}-f_\text{éloi}$. La moyenne des deux valeurs
extrêmes donne-t-elle $f_\mathrm{E}$ ? Comment retrouver $v_\mathrm{E}$ et $f_\mathrm{E}$ à partir des fréquences extrêmes ?
\item La fréquence varie surtout pendant que la voiture parcourt les quelques dizaines de mètres les plus proches de
l'auditeur. Estimer la durée de cette transition et comparer les situations $b=\qty{10}{m}$ et $b=\qty{50}{m}$.
\end{questions}
\end{blocEdition}

\editionPage{defi-ex-24}{D24 · Où placer l'œil ? Le cercle oculaire}{Recherche et rédaction}{Avancé · 45 à 60~min · Chapitre 24\hfill\hyperref[defi-cor-24-1]{Corrigé p.~\pageref*{defi-cor-24-1}}}
\begin{blocEdition}{defi-ex-24}
Toute la lumière collectée par l'objectif d'une lunette ressort en passant par un petit disque situé derrière
l'oculaire : le \textbf{cercle oculaire}, image de l'objectif donnée par l'oculaire. C'est là que l'observateur doit placer
la pupille de son œil. On étudie une lunette 70/700 (objectif de diamètre $D=\qty{70}{mm}$ et de distance focale
$f'_1=\qty{700}{mm}$) utilisée avec plusieurs oculaires de distance focale $f'_2$. On rappelle la relation de conjugaison d'une
lentille mince, $\dfrac{1}{\overline{OA'}}-\dfrac{1}{\overline{OA}}=\dfrac{1}{f'}$, et l'expression du grandissement, $\gamma=\dfrac{\overline{OA'}}{\overline{OA}}$. Le
programme calcule, pour chaque oculaire, le grossissement, la distance $x$ entre l'oculaire et le cercle oculaire, et le
diamètre de ce cercle, en millimètres.
\begin{codecote}[0.7]
\begin{pythoncol}
D, f1 = 70, 700                  # mm
for f2 in [40, 25, 10, 4]:       # mm, oculaires
    G = f1 / f2
    x = f2 * (f1 + f2) / f1      # cercle oculaire
    print(f2, round(G, 1), round(x, 1), round(D/G, 1))
\end{pythoncol}
\cote
\begin{consolecol}
40 17.5 42.3 4.0
25 28.0 25.9 2.5
10 70.0 10.1 1.0
4 175.0 4.0 0.4
\end{consolecol}
\end{codecote}
\par\smallskip
\begin{questions}
\item L'objectif se trouve à la distance $f'_1+f'_2$ en avant de l'oculaire. Montrer que son image par l'oculaire est réelle
et située à la distance $x=\dfrac{f'_2\,(f'_1+f'_2)}{f'_1}$ derrière l'oculaire. Calculer $x$ pour l'oculaire de \qty{10}{mm}.
\item Montrer que le diamètre du cercle oculaire vaut $d'=\dfrac{D\,f'_2}{f'_1}=\dfrac{D}{G}$.
\item Interpréter la sortie du programme. Pourquoi les oculaires de courte focale obligent-ils à approcher l'œil très
près du verre ?
\item La nuit, la pupille de l'œil s'ouvre jusqu'à \qty{7}{mm} environ. Quel grossissement minimal faut-il pour que toute la
lumière collectée par l'objectif entre dans l'œil ? Quel est l'oculaire de plus grande focale utilisable ?
\item Pour un astre étendu comme une planète, l'éclat de l'image est proportionnel à $d'^2$. Comparer l'éclat obtenu avec
les oculaires de \qty{10}{mm} et de \qty{40}{mm}, puis relier la règle du grossissement maximal utile ($G\approx2D$, $D$ en mm) au
diamètre du cercle oculaire.
\end{questions}
\end{blocEdition}

\editionPage{defi-ex-25}{D25 · Mesurer la constante de Planck avec des LED}{Recherche et rédaction}{Avancé · 45 à 60~min · Chapitre 25\hfill\hyperref[defi-cor-25-1]{Corrigé p.~\pageref*{defi-cor-25-1}}}
\begin{blocEdition}{defi-ex-25}
Une diode électroluminescente (LED) n'émet de la lumière que si la tension à ses bornes dépasse une tension de seuil
$U_\mathrm{s}$. À ce seuil, chaque électron qui traverse la diode reçoit l'énergie $e\,U_\mathrm{s}$, qu'il cède en émettant un photon
d'énergie voisine de $h\,c/\lambda$. Un élève mesure la tension de seuil de cinq LED.
\begin{center}
\begin{tabular}{lccccc}
\toprule
LED & infrarouge & rouge & jaune & verte & bleue \\
$\lambda$ (nm) & 940 & 630 & 590 & 525 & 470 \\
$U_\mathrm{s}$ (V) & 1,28 & 1,91 & 2,04 & 2,29 & 2,56 \\
\bottomrule
\end{tabular}
\end{center}
Le programme effectue la régression linéaire de $U_\mathrm{s}$ en fonction de $1/\lambda$ ($e=\qty{1,602e-19}{C}$ ;
$c=\qty{2,998e8}{m.s^{-1}}$).
\begin{python}
lam = [940, 630, 590, 525, 470]          # nm
Us = [1.28, 1.91, 2.04, 2.29, 2.56]      # V
x = [1e9 / l for l in lam]               # 1/lambda (m^-1)
n = len(x)
mx, my = sum(x) / n, sum(Us) / n
num = sum((xi - mx) * (yi - my) for xi, yi in zip(x, Us))
den = sum((xi - mx)**2 for xi in x)
a = num / den                            # coefficient directeur (V.m)
h = a * 1.602e-19 / 2.998e8
print(round(a * 1e6, 3), round(h * 1e34, 2))
\end{python}
\begin{console}
1.203 6.43
\end{console}
\begin{questions}
\item Montrer que, si $e\,U_\mathrm{s}=\dfrac{h\,c}{\lambda}$, la tension de seuil est proportionnelle à $1/\lambda$, avec un coefficient
directeur égal à $\dfrac{h\,c}{e}$.
\item Que contient la liste \code{x} ? Que calcule la variable \code{a} ? Dans quelle unité le programme affiche-t-il sa
valeur ?
\item En déduire la valeur de $h$ obtenue par l'élève et l'écart relatif avec la valeur de référence,
$h=\qty{6,626e-34}{J.s}$.
\item Pour chaque LED, comparer $U_\mathrm{s}$ à $\dfrac{h\,c}{e\lambda}$. L'écart est-il aléatoire ou systématique ? Proposer des causes.
\item L'élève a ajouté une LED infrarouge, dont la lumière est invisible. Quel est l'intérêt de cette mesure ? Comment
repérer que cette LED s'allume ?
\end{questions}
\end{blocEdition}

\editionPage{defi-ex-26}{D26 · Où passe la moitié de l'énergie ?}{Recherche et rédaction}{Très exigeant · 50 à 60~min · Chapitre 26\hfill\hyperref[defi-cor-26-1]{Corrigé p.~\pageref*{defi-cor-26-1}}}
\begin{blocEdition}{defi-ex-26}
Un condensateur de capacité $C=\qty{1,0}{mF}$, initialement déchargé, est chargé à travers une résistance $R$ par un
générateur idéal de tension $E=\qty{10}{V}$. Un condensateur chargé sous la tension $u$ stocke l'énergie $\frac12\,C\,u^2$. Le
programme calcule, par une somme de petits termes $R\,i^2\,\mathrm{d}t$, l'énergie $W_\mathrm{J}$ dissipée par effet Joule dans la
résistance pendant la charge, pour trois valeurs de $R$, et la compare à l'énergie finalement stockée.
\begin{codecote}[0.66]
\begin{pythoncol}
from math import exp
E, C = 10.0, 1.0e-3               # V, F
for R in [10, 100, 1000]:         # ohms
    tau = R * C
    dt = tau / 1000
    W, t = 0.0, 0.0
    while t < 10 * tau:
        i = E / R * exp(-t / tau)
        W = W + R * i**2 * dt
        t = t + dt
    print(R, round(W, 3), 0.5 * C * E**2)
\end{pythoncol}
\cote
\begin{consolecol}
10 0.05 0.05
100 0.05 0.05
1000 0.05 0.05
\end{consolecol}
\end{codecote}
\par\smallskip
\begin{questions}
\item Calculer la charge finale du condensateur. Le générateur fait circuler cette charge sous la tension constante $E$ :
montrer qu'il fournit l'énergie $W_\mathrm{g}=C\,E^2$.
\item Calculer l'énergie stockée en fin de charge. Quelle part de l'énergie fournie a été perdue ? Quel est le rendement de
la charge ?
\item Sachant que $i(t)=\dfrac{E}{R}\,\mathrm{e}^{-t/\tau}$, montrer par le calcul d'une intégrale que $W_\mathrm{J}=\displaystyle\int_0^{+\infty}R\,i^2\,\mathrm{d}t=\frac12\,C\,E^2$.
\item Interpréter la sortie du programme. Pourquoi peut-on arrêter la somme à $t=10\tau$ ?
\item On charge maintenant le condensateur en deux étapes : d'abord sous $E/2$ jusqu'à charge complète, puis sous $E$.
Calculer l'énergie dissipée et le rendement. Que se passerait-il avec $N$ étapes ?
\end{questions}
\end{blocEdition}

\editionPage{defi-ex-27}{D27 · Mesurer plus pour mesurer mieux ?}{Recherche et rédaction}{Très exigeant · 50 à 60~min · Chapitre 27\hfill\hyperref[defi-cor-27-1]{Corrigé p.~\pageref*{defi-cor-27-1}}}
\begin{blocEdition}{defi-ex-27}
Des élèves mesurent la durée de chute $t$ d'une bille lâchée sans vitesse initiale d'une hauteur $h$, pour en déduire
$g=\dfrac{2h}{t^2}$. Leurs mesures sont dispersées avec un écart-type $s=\qty{0,013}{s}$. Pour des grandeurs indépendantes,
$u(x+y)^2=u(x)^2+u(y)^2$ et, pour une constante $k$, $u(k\,x)=|k|\,u(x)$. Le programme simule \num{2000} séries de $N=40$ mesures.
\begin{codecote}[0.62]
\begin{pythoncol}
import random
from statistics import mean, stdev
s, N = 0.013, 40
moyennes = []
for k in range(2000):
    serie = [random.gauss(0.452, s)
             for j in range(N)]
    moyennes.append(mean(serie))
print(round(stdev(moyennes), 5),
      round(s / N**0.5, 5))
\end{pythoncol}
\cote
\begin{consolecol}
0.00205 0.00206
\end{consolecol}
\par\medskip
\fontsize{9.5}{12.5}\selectfont
Durées moyennes mesurées (40 lâchers par hauteur) :\par\smallskip
\begin{tabular}{@{}lc@{}}
\toprule
$h$ (m) & $t_\mathrm{mes}$ (s) \\
\midrule
0,20 & 0,2119 \\
0,40 & 0,2956 \\
0,60 & 0,3597 \\
0,80 & 0,4139 \\
1,00 & 0,4615 \\
\bottomrule
\end{tabular}
\end{codecote}
\par\smallskip
\begin{questions}
\item Montrer que la moyenne de $N$ mesures indépendantes, d'écart-type $s$, a pour incertitude-type $\dfrac{s}{\sqrt N}$.
\item Que calcule le programme ? Interpréter les deux valeurs affichées.
\item Le dispositif de lâcher ajoute en réalité un retard $\delta=\qty{0,010}{s}$ à chaque durée. Pour $h=\qty{1,000}{m}$ et
$g=\qty{9,81}{m.s^{-2}}$, calculer la durée de chute réelle. Exprimer l'écart normalisé entre la moyenne de $N$ mesures et
cette durée. À partir de combien de mesures le retard est-il détecté ?
\item Calculer la valeur de $g$ déduite de la moyenne de 40 mesures, son incertitude-type, avec $u(h)=\qty{0,0020}{m}$ et
$\dfrac{u(g)}{g}=\sqrt{\left(\dfrac{u(h)}{h}\right)^2+\left(2\,\dfrac{u(t)}{t}\right)^2}$, puis l'écart normalisé avec \qty{9,81}{m.s^{-2}}. Mesurer plus a-t-il permis de
mesurer mieux ?
\item Montrer que $t_\mathrm{mes}=\sqrt{\dfrac{2}{g}}\,\sqrt h+\delta$. Exploiter le tableau pour déterminer $\delta$ et $g$.
\end{questions}
\end{blocEdition}
